\begin{claim}{F.33}{Proved}{Exact paired restriction kernel for opposite first characters}
Use the actual bases $Y_i^-=X_i$ and
$Y_i^+=\mathcal B_+(X_i)$ of F.23.13.
Let $k$ be a finite field of odd characteristic,
and let $\chi,\chi':V_\pm\to k^\times$ be the
continuous characters with values
\[
 \begin{aligned}
 (\chi(Y_i^\pm))_{i=1}^7&=(r,s,t,u,v,w,z),\\
 (\chi'(Y_i^\pm))_{i=1}^7&=(-r,s,t,u,v,w,z),\qquad
             \langle u,v,w,z\rangle\ne\{1\}.
 \end{aligned}                                                \tag{F.33.1}
\]
All seven displayed values are nonzero; $r^2=1$
is included. The two characters have the same
restriction $\eta$ to the actual subgroup
$\mathcal K_\pm$. The summed continuous cocycle
restriction has the exact kernel
\[
 \begin{split}
 \ker\bigl[Z^1_\chi(V_\pm,k)\oplus
 Z^1_{\chi'}(V_\pm,k)
   &\longrightarrow Z^1_\eta(\mathcal K_\pm,k)\bigr]\\
   &=\{(b(\chi-1),-b(\chi'-1)):b\in k\}.
 \end{split}                                                  \tag{F.33.2}
\]
It has dimension one, and the restriction
map has rank $13$ on the original $14$
generator coordinates. The induced summed
$H^1$ restriction is injective.

More generally, the explicit $64$ by $13$
Laurent matrix used below has rank $13$
over every field, in every characteristic,
for every nonzero tuple $(R,s,t,u,v,w,z)$
with $(u,v,w,z)\ne(1,1,1,1)$.
\end{claim}

\begin{proof}
\emph{The original subgroup and source words.}
First work with
$K=\mathcal K_-\le R_-\le V_-=V$.
The original retained quotient
$\pi:V\to\widehat F(X_1,X_2,X_3)$ kills
$X_4,\ldots,X_7$, and $R_-=\ker\pi$.
The character $\chi'/\chi$ is exactly
$\epsilon\pi$, where
$\epsilon(X_1)=-1$ and
$\epsilon(X_2)=\epsilon(X_3)=1$.
Thus $\chi|R_-=\chi'|R_-$, proving their
agreement on $K$ by this actual quotient map.
The normal-generation equality of F.21,
transported through the original marginal
map of F.22 and then through the abelian
scalar map, gives
\[
       \chi(K)=\chi(R_-)=\langle u,v,w,z\rangle.
                                                               \tag{F.33.3}
\]
Indeed conjugation does not change an
abelian scalar image, and the finite image
of $K$ is already closed. In particular
$\eta\ne1$.

Retain the eight actual negative collision
seeds in the following order:
\[
 \begin{gathered}
 \omega_0,\ldots,\omega_7:\quad
 (0,-1,1),(0,-1,-1),(2,-1,1),(2,-1,-1),\\
 (3,-1,1),(3,-1,-1),(4,-1,1),(4,-1,-1).
 \end{gathered}                                                \tag{F.33.4}
\]
The triples are the original F.23 labels
$\mathrm{cross}_{e,-1,\varepsilon}$.
Their scalar characters are respectively
$u,u,v,v,w,w,z,z$.
Write $\beta_j$ for the actual negative
old-source automorphism with zero-based
index $j$, including its original
decoration and orientation. The $64$ words
used in the proof are, in order,
\[
 \begin{gathered}
 \omega_0,\ldots,\omega_7,\qquad
 (\beta_j(\omega_0),\ldots,\beta_j(\omega_7))
       \quad\text{successively for}\\
                         j=2,9,8,10,1,0,7.
 \end{gathered}                                                \tag{F.33.5}
\]
These are the original F.24 source rows.
Every displayed word belongs to $K$:
the seeds lie in $K$, and each original
old-source conjugation preserves $K$.
No larger braid group or replacement
source map is used. Since the old sources
are pure, their action preserves every
abelian character, so every word in
(F.33.5) retains its seed character.

\emph{The coefficient ring and the principal-line quotient.}
Set
\[
 \begin{gathered}
 A=\mathbb Z[R^{\pm1},s^{\pm1},t^{\pm1},u^{\pm1},
                         v^{\pm1},w^{\pm1},z^{\pm1}],\\
                         B=A[\xi]/(\xi^2-R).
 \end{gathered}
                                                               \tag{F.33.6}
\]
The element $\xi$ is a unit, with
$\xi^{-1}=\xi/R$. The map from the
original seven-variable Laurent ring
to $B$ sends its variables to
$(\xi,s,t,u,v,w,z)$. Its inverse sends
$\xi$ to the original first variable
and $R$ to its square, fixing the other
six variables. Both composites fix all
generators. The monic quadratic relation
gives the free $A$-basis $1,\xi$;
explicitly, for every integer $e$,
\[
       \xi^e=R^{(e-\bar e)/2}\xi^{\bar e},
       \qquad \bar e\in\{0,1\},\quad
                         \bar e\equiv e\pmod2.
                                                               \tag{F.33.7}
\]
This includes negative exponents and
defines the exact coefficient map.

For each word in (F.33.5), write its
original universal derivation row
uniquely as $a+\xi b$, with $a,b\in A^7$.
Let $E$ be the $64$ by $14$ matrix whose
row is $(b,a)$, in that order.
The exact Fox fundamental identity for
the word, whose character has no $\xi$
component, gives
\[
 -b_1+a_1+\sum_{j=2}^7b_j(\theta_j-1)=0,
       \qquad(\theta_2,\ldots,\theta_7)=(s,t,u,v,w,z).
                                                               \tag{F.33.8}
\]
Consequently $Ep=0$ for the actual vector
\[
 p=(-1,s-1,t-1,u-1,v-1,w-1,z-1;\,1,0,0,0,0,0,0)^T.
                                                               \tag{F.33.9}
\]

Here is an explicit quotient with its
inverse, over $A$ itself. For
$(x_1,\ldots,x_7;y_1,\ldots,y_7)$ put
\[
 \begin{aligned}
 \lambda&=y_1,&q_1&=x_1+\lambda,\\
 q_j&=x_j-(\theta_j-1)\lambda\quad(2\le j\le7),&
 q_{6+j}&=y_j\quad(2\le j\le7).
 \end{aligned}                                                \tag{F.33.10}
\]
The target coordinates are $(q_1,\ldots,q_{13};\lambda)$.
The inverse is
\[
 \begin{aligned}
 x_1&=q_1-\lambda,&y_1&=\lambda,\\
 x_j&=q_j+(\theta_j-1)\lambda\quad(2\le j\le7),&
 y_j&=q_{6+j}\quad(2\le j\le7).
 \end{aligned}                                                \tag{F.33.11}
\]
Substitution proves both inverse
identities coordinate by coordinate.
This subtracts $\lambda p$ and records
the unique removed parameter. By
(F.33.8), $E(x,y)=Nq$, where the row of
the $64$ by $13$ matrix $N$ is
\[
              (b_1,\ldots,b_7,a_2,\ldots,a_7).
                                                               \tag{F.33.12}
\]
There is no division in this quotient.

\emph{Four exhaustive exact row certificates.}
The field specializations are covered by
the four disjoint charts
\[
 \begin{array}{c|l|c}
 \mathsf U&u\ne1&\delta=u-1\\
 \mathsf V&u=1,\ v\ne1&\delta=v-1\\
 \mathsf W&u=v=1,\ w\ne1&\delta=w-1\\
 \mathsf Z&u=v=w=1,\ z\ne1&\delta=z-1.
 \end{array}                                                  \tag{F.33.13}
\]
Setting the preceding variables to one
is the displayed Laurent ring
homomorphism for that chart.
Every row operation is either multiplication
by a Laurent monomial, division by the
active $\delta$ after exact divisibility
has been checked in every entry, or
the cleared elimination
$L\mapsto pL-cP$, where $p$ is the
pivot entry of $P$ and $c$ is the
corresponding entry of $L$.
The inverse to each divided normalization
is multiplication by its recorded
monomial inverse and power of $\delta$.
Thus all resulting rows lie in the
specialized original row space.

The producer permits tests for factors
$R-1$ and $2$ during its search, but
every recorded division count for those
factors is zero, in all operations and
all unit witnesses. They are not inverted
in this proof. The only actual divisions
are by character monomials and the
nonzero active difference in (F.33.13).
The respective operation counts are
$674,483,374,232$; the respective
active-difference division counts are
$17,15,15,9$.

The following table prints every single
pivot as $(\text{row},\text{column};\text{coefficient})$.
Both row and column indices are zero-based,
and the column order is exactly (F.33.12).
\[
\begin{array}{c|c|c|c|c}
 &\mathsf U&\mathsf V&\mathsf W&\mathsf Z\\\hline
1&(11,4;stu)&(0,3;1)&(0,3;1)&(0,3;1)\\
2&(8,9;1)&(1,9;1)&(1,9;1)&(1,9;1)\\
3&(13,5;stu)&(3,10;s)&(2,4;1)&(2,4;1)\\
4&(14,12;1)&(11,4;st)&(3,10;1)&(3,10;1)\\
5&(0,3;1)&(13,5;st)&(13,5;st)&(4,5;1)\\
6&(3,10;su)&(14,12;1)&(14,12;1)&(5,11;1)\\
7&(16,8;s)&(18,7;-1)&(20,7;-1)&(14,12;1)\\
8&(6,6;1)&(6,6;-1)&(4,0;-stz)&(39,7;-1)\\
9&(40,0;1)&(5,11;1)&(28,8;1)&(62,0;st)\\
10&(1,2;s)&(10,0;-st)&(6,6;st)&(7,8;-1)\\
11&(48,7;1)&(34,8;-1)&(5,11;1)&(6,6;st)\\
12&(5,11;1)&\text{block below}&(29,2;s)&(47,2;s)\\
13&(49,1;1)&\text{block below}&(52,1;-1)&(55,1;1).
\end{array}                                                   \tag{F.33.14}
\]
Every entry is a signed Laurent monomial,
and every pivot row is zero in its
preceding pivot columns. The remaining
two columns of chart $\mathsf V$ are
$1,2$. In the actual recorded order
of rows $35,51$, their exact block is
\[
 M=\begin{pmatrix}
 -(t-1)(st+1)&s(st-t+1)\\
 st-s+1&-s(s-1)
 \end{pmatrix},\qquad \det M=-s^2t.
                                                               \tag{F.33.15}
\]
These two full rows vanish in all eleven
previous pivot columns. Direct multiplication
gives $M\operatorname{adj}(M)=
\operatorname{adj}(M)M=-s^2t I_2$.
Thus the adjugate gives the two remaining
coordinate rows with unit coefficient
$-s^2t$. In the reversed row order the
block is $SM$, where
$S=\left(\begin{smallmatrix}0&1\\1&0\end{smallmatrix}\right)$;
$S^2=I_2$ and $\det(SM)=s^2t$.
This is the exact comparison between
the two row-order descriptions.

The fifty single pivots and this one
two-column block prove rank $13$ in
all four charts. The proof is over
the integers, with only the stated
nonzero field divisions. It therefore
establishes the matrix assertion in
every characteristic, including $2$,
and for every $R\ne0$, including $R=1$.

\emph{The two characters and their exact kernel.}
Now specialize $R=r^2$ in the given odd
characteristic field. A derivation row
$a+\xi b$ becomes $a+rb$ for $\chi$
and $a-rb$ for $\chi'$.
For their generator-value columns $c,c'$,
the actual summed evaluation is
\[
 \begin{gathered}
 (a+rb)c+(a-rb)c'=bx+ay,\\
 \begin{aligned}
 x&=r(c-c'),&y&=c+c',\\
 c&=(y+x/r)/2,&c'&=(y-x/r)/2.
 \end{aligned}
 \end{gathered}                                                \tag{F.33.16}
\]
The displayed two maps are inverse;
their only divisors are the nonzero
elements $2,r$. Together with
(F.33.10)--(F.33.12), rank $13$ gives
$\ker E=kp$. The inverse map in
(F.33.16) sends $\lambda p$ to
\[
 \left(\frac{\lambda}{2r}(\chi(X_i)-1)_{i=1}^7,
       -\frac{\lambda}{2r}(\chi'(X_i)-1)_{i=1}^7\right).
                                                               \tag{F.33.17}
\]
Thus vanishing on the actual $64$ words
already forces exactly the generator
values in (F.33.2).

Every assignment of the seven generator
values extends uniquely to a continuous
crossed homomorphism with the fixed
scalar character: it is the unique
homomorphism from the original free
profinite group to the finite group
$\operatorname{Aff}_1(k)$ with those
generator images. Hence the generator
identity (F.33.17) is an identity of
continuous cocycles on $V$.
Conversely the two principal cocycles
in (F.33.2) sum to zero on all of $K$,
because $\chi|K=\chi'|K$.
This proves both kernel inclusions.
The vector in (F.33.2) is nonzero for
$b\ne0$, since $\eta\ne1$; its dimension
is one and the rank is $14-1=13$.

For the cohomology statement, the two
principal cocycles restrict to principal
$\eta$-cocycles, so the summed map is
well-defined on $H^1$. A kernel class
has representatives whose sum on $K$
is $b(\eta-1)$. Subtract $b(\chi-1)$
from the first representative. The new
representatives have zero summed
restriction, so (F.33.2) makes both
principal on $V$. Their original
cohomology classes are therefore zero.

\emph{Literal-word audit and the positive marginal.}
The independent word verifier imports
no producer arithmetic and no symbolic
algebra package. It uses generic
crossed-homomorphism multiplication
$(\alpha,D)(\beta,E)=(\alpha\beta,D+\alpha E)$
and inverse $(\alpha^{-1},-\alpha^{-1}D)$
to reconstruct the original nineteen
negative F.23 braid words.
At a positive crossing it applies
\[
 (x_j,x_{j+1};c_j,c_{j+1})\longmapsto
 (x_jx_{j+1}x_j^{-1},x_j;x_jc_{j+1},c_j),
                                                               \tag{F.33.18}
\]
and at a negative crossing it applies
$(x_{j+1},x_{j+1}^{-1}x_jx_{j+1};
c_{j+1},x_{j+1}^{-1}c_j)$.
The eighth meridian is the identity;
the actual decoration is $c_8^{-1}$.
The verifier checks every pure label
permutation, decoration character and
derivative, all eleven full seven-coordinate
actions, all $64$ original rows, the
principal identities, and every recorded
normalization, elimination, pivot and
adjugate. Its $252309$ exact checks pass.
The separate root replay has $32427$
checks, including the full inverse
coordinate maps, and the separate
actual-divisor audit checks all $51$
unit witnesses. The evidence manifest
binds these inputs, programs and outputs.

For the positive marginal retain the
actual based maps of F.23.13,
\[
 \mathcal B_+(g)=w_+\Pi_+(g)w_+^{-1},\qquad
 \mathcal B_+^{-1}(g)=\Pi_+^{-1}(w_+^{-1}gw_+).
                                                               \tag{F.33.19}
\]
Their composites cancel to the identity.
The original filling squares identify
their restrictions $V_-\leftrightarrow V_+$
and $R_-\leftrightarrow R_+$, and the
actual source comparison gives
$\mathcal B_+(\mathcal K_-)=\mathcal K_+$.
Precomposition sends a positive character
and cocycle to
$(\chi\mathcal B_+,d\mathcal B_+)$;
precomposition with $\mathcal B_+^{-1}$
is its inverse. Since $Y_i^+=\mathcal B_+(X_i)$,
this preserves all seven displayed
values, the paired coordinate maps,
and the actual restriction sums.
It also sends every negative word in
(F.33.5) to its exact positive image.
The equality of cocycles, both kernel
inclusions and the cohomology injection
therefore transport to the positive
marginal. All groups and coefficient
maps used here are continuous, with
finite coefficient targets throughout.
This proves the claim.
\end{proof}
